\section{preempts}

In this section, we shall discuss the range and responses to opening preempts or direct jump overcalls.
Here are two main "dimensions" to properly describe a preemptive hand:

\begin{itemize}
    \setlength{\itemsep}{0pt}
    \item defensiveness: the expected tricks when defending, for example, general HCP
    \item offensiveness: the expected tricks when preempting, for example, shape and KQJT in long suit
\end{itemize}

In general, a preemptive range can be defined as a subset on the defensiveness-offensiveness plane.

\begin{itemize}
    \setlength{\itemsep}{0pt}
    \item defensiveness: the expected tricks when defending, for example, general HCP
    \item offensiveness: the expected tricks when preempting, for example, shape and KQJT in long suit
\end{itemize}

\subsection{style}

After some advices and top-player's play, I think we'll use the following approach:

\begin{itemize}
    \setlength{\itemsep}{0pt}
    \item baseline: preempts follows 2-3-4 rule and shows 6+ cards. 
    \item seat adjustments: 
    \begin{itemize}
        \setlength{\itemsep}{0pt}
        \item opening soundness: 2nd > 1st > 3rd
        \item jump overcall soundness: 2nd > 1st > 3rd
    \end{itemize}
    \item Vul adjustments: 
    \begin{itemize}
        \setlength{\itemsep}{0pt}
        \item V/NV: 3X always promises 7-cards
        \item V/V: 3X usually promises 7-cards
    \end{itemize}
    \item 3+ level can be lighter with less defensiveness or more offensiveness. but 2-level usually won't unless partner PH.
    \item We compensate this with a lighter and wider overcall range.
\end{itemize}

\subsection{2-3-4 rule revisited}

The 2-3-4 rule is a good starting point to discuss. I suggest the following modification according to vulnerability and seats:

\begin{center}
\begin{NiceTabular}{lcccc}[hvlines]
 & V/NV & V/V & NV/NV & NV/V \\
1st open & 2 & 3 & 3 & 4 \\
2nd open & 2 & 3- & 3 & 4- \\
3rd open & 2+ & 3+ & \cellcolor{lime}- & \cellcolor{lime}- \\
overcall & 2+ & 3 & 3 & 4 \\
overcall (partner PH) & 2+ & 3+ & \cellcolor{tiffanyblue}- & \cellcolor{tiffanyblue}- \\
\end{NiceTabular}
\end{center}

The color indicates situations where we might preempt with 5-cards:
\begin{itemize}
    \setlength{\itemsep}{0pt}
    \item {\color{green}green} NV 3-rd seat opening will be frequent 5-card.
    \item {\color{cyan}blue} occasionally random weak 5-card (especially 2S). but partner usually treat as a normal 6-card preempt.
\end{itemize}

\subsection{opening pre}

We use multi + Muiderberg, which is consistent with our defense to 1N. The number of cards promised for multi follows the above table, 
and the number of minor cards promised for Muiderberg is always 5-5.

\bin{\currfiledir/seq/pre}

\subsection{opening pre interfered}

\bin{\currfiledir/seq/pre_interfered}

\subsection{after partner's pre vs 1X}

We simply plays natrual (2N = nat inv, cue = inv+ raise) responses.

% \bin{\currfiledir/seq/resp}

\subsection{notes}

\bin{\currfiledir/seq/note}

\begin{itemize}
    \setlength{\itemsep}{0pt}
    \item jump to 4M = s/o except after multi 2D - 4H
    \item After preempt, sometimes psychic new suit or 3N, or just lead-directing.
    \item 4M opening may be made with 5 or fewer Losers (especially in 3rd seat). In this case, the opener will double if interfered to show "my 4M is to make".
\end{itemize}
